\documentclass[17pt,compress]{beamer}
\usepackage{beamerthemesplit}
\mode<presentation>
{
  \usetheme{Warsaw}
  \setbeamercovered{transparent}
  \setbeamertemplate{navigation symbols}{}
}
% Taken from Fernando's slides.
\usepackage{ae,aecompl}
\usepackage[scaled=.95]{helvet}

\usepackage[english]{babel}
\usepackage[latin1]{inputenc}
\usepackage[T1]{fontenc}

\newcounter{saveenumi}
\newcommand{\seti}{\setcounter{saveenumi}{\value{enumi}}}
\newcommand{\conti}{\setcounter{enumi}{\value{saveenumi}}}

\definecolor{blue}{rgb}{0.16,0.32,0.75}
\setbeamercolor{structure}{fg=blue}
\author[FOSSEE]{}
\institute[IIT Bombay]{}
\date[]{}
% \setbeamercovered{transparent}

% theme split
\usepackage{verbatim}
\newenvironment{colorverbatim}[1][]%
{%
\color{blue}
\verbatim
}%
{%
\endverbatim
}%

\usepackage{mathpazo,courier,euler}
\usepackage{listings}
\lstset{language=sh,
    basicstyle=\ttfamily\bfseries,
  showstringspaces=false,
  keywordstyle=\color{black}\bfseries}

% logo
\logo{\includegraphics[height=1.30 cm]{3t-logo.pdf}}
\logo{\includegraphics[height=1.30 cm]{fossee-logo.png}

\hspace{7.5cm}
\includegraphics[scale=0.3]{fossee-logo.png}\\
\hspace{281pt}
\includegraphics[scale=0.80]{3t-logo.pdf}}

\begin{document}

\sffamily \bfseries
\title
[Saving Plots]
{Python - Saving Plots}
\author
[FOSSEE, IIT - Bombay]
{\small Talk to a Teacher\\{\color{blue}\url{http://spoken-tutorial.org}}\\National Mission on Education
 through ICT\\{\color{blue}\url{http://sakshat.ac.in}} \\[0.5cm]{\tiny Script by: Hardik ghaghada \\ Naration by: Hardik Ghaghada \\ 13 August 2015}}

% slide 1
\begin{frame}
   \titlepage
\end{frame}
%%%%%%%%%%%%%%%%%%%%%%%%%%%%%%%%%%%%%%%%%%%%%%%%%%%%%%%%%%%%%%%%%%%%%%%%%%%%%%%%
\begin{frame}
\frametitle{Objectives} 
\label{sec-2}
In this tutorial, we will learn -\pause
\begin{itemize}
\item Save plots using \texttt{savefig()}function. \pause
\item Save plots in different formats.
\end{itemize}
\end{frame}
%%%%%%%%%%%%%%%%%%%%%%%%%%%%%%%%%%%%%%%%%%%%%%%%%%%%%%%%%%%%%%%%%%%%%%%%%%%%%%%%
\begin{frame}
\frametitle{System Specifications}\pause
\begin{itemize}
\item Slackware Linux 14.1\pause
\item \texttt{Python 2.7.5} \pause
\item \texttt{IPython 2.3.0}
\end{itemize}
\end{frame}
%%%%%%%%%%%%%%%%%%%%%%%%%%%%%%%%%%%%%%%%%%%%%%%%%%%%%%%%%%%%%%%%%%%%%%%%%%%%%%%%
\begin{frame}
\frametitle{Pre-requisites}
\label{sec-3}
Spoken tutorial on -
\begin{itemize}
\item ``Using plot interactively''
\end{itemize}
\end{frame}
%%%%%%%%%%%%%%%%%%%%%%%%%%%%%%%%%%%%%%%%%%%%%%%%%%%%%%%%%%%%%%%%%%%%%%%%%%%%%%%%
\begin{frame}[fragile]
\frametitle{Creating a basic plot}
\label{sec-4}
Plot a sine wave from -3pi to 3pi.\pause
\begin{small}
\lstset{language=Python}
\begin{lstlisting}
In []: x = linspace(-3*pi,3*pi,100)
In []: plot(x, sin(x))
\end{lstlisting}
\end{small}
\end{frame}
%%%%%%%%%%%%%%%%%%%%%%%%%%%%%%%%%%%%%%%%%%%%%%%%%%%%%%%%%%%%%%%%%%%%%%%%%%%%%%%%
\begin{frame}[fragile]
\frametitle{savefig()}
\label{sec-5}
\begin{itemize}
\item \texttt{savefig()} - to save plots\pause
\label{sec-5_1}%

Syntax: \texttt{savefig(fname)}\pause

\item example
\label{sec-5_2}%
\begin{itemize}
\item savefig(`/home/fossee/sine.png')\\\pause
\label{sec-5_2_1}%
\begin{itemize}
\item file sine.png saved to the folder /home/fossee\pause
\item .png - file type
\end{itemize}
\end{itemize} % ends low level
\end{itemize} % ends low level
\end{frame}
%%%%%%%%%%%%%%%%%%%%%%%%%%%%%%%%%%%%%%%%%%%%%%%%%%%%%%%%%%%%%%%%%%%%%%%%%%%%%%%%
\begin{frame}
\frametitle{More on savefig()}
\label{sec-6}
\begin{itemize}
\item Recall\\
\label{sec-6_1}%
\begin{itemize}
\item .png - file type\pause
\end{itemize}
\item File types supported\pause
\label{sec-6_2}%
\begin{itemize}

\begin{footnotesize}
\item .pdf - PDF(Portable Document Format)\pause\\
\label{sec-6_2_1}%
\item .ps - PS(Post Script)\pause\\
\label{sec-6_2_2}%
\item .eps - Encapsulated Post Script\\
\label{sec-6_2_3}%
\verb~to be used with~ \LaTeX{} \verb~documents~\pause
\item .svg - Scalable Vector Graphics\\
\label{sec-6_2_4}%
\verb~vector graphics~\pause
\item .png - Portable Network Graphics\\
\label{sec-6_2_5}%
\verb~supports transparency~
\end{footnotesize}
\end{itemize} % ends low level
\end{itemize} % ends low level
\end{frame}
%%%%%%%%%%%%%%%%%%%%%%%%%%%%%%%%%%%%%%%%%%%%%%%%%%%%%%%%%%%%%%%%%%%%%%%%%%%%%%%%
\begin{frame}
\frametitle{Exercise 1}
\label{sec-7}
Save the sine plot in the format EPS which can be embedded in \LaTeX{} documents.
\end{frame}
%%%%%%%%%%%%%%%%%%%%%%%%%%%%%%%%%%%%%%%%%%%%%%%%%%%%%%%%%%%%%%%%%%%%%%%%%%%%%%%%
\begin{frame}
\frametitle{Exercise 2}
\label{sec-8}
Save the sine plot in PDF, PS and SVG formats.
\end{frame}
%%%%%%%%%%%%%%%%%%%%%%%%%%%%%%%%%%%%%%%%%%%%%%%%%%%%%%%%%%%%%%%%%%%%%%%%%%%%%%%%
\begin{frame}
\frametitle{Summary}
\label{sec-9}
In this tutorial, we have learnt to -
\begin{itemize}
\item Save plots using the \texttt{savefig()} function.
\item Save the plots in different formats.
\begin{itemize}
\item pdf, ps, png, svg, eps
\end{itemize}
\item Locate files in the file system.
\end{itemize}
\end{frame}
%%%%%%%%%%%%%%%%%%%%%%%%%%%%%%%%%%%%%%%%%%%%%%%%%%%%%%%%%%%%%%%%%%%%%%%%%%%%%%%%
\begin{frame}
\frametitle{Evaluation}
\label{sec-10}
\begin{enumerate}
\item Which command is used to save a plot.\pause
	\begin{itemize}
	\item \texttt{saveplot()}
	\item \texttt{savefig()}
	\item \texttt{savefigure()}
	\item \texttt{saveplt()}
	\end{itemize}
	\seti
\end{enumerate}
\end{frame}
%%%%%%%%%%%%%%%%%%%%%%%%%%%%%%%%%%%%%%%%%%%%%%%%%%%%%%%%%%%%%%%%%%%%%%%%%%%%%%%%
\begin{frame}
\frametitle{Evaluation contd..}
\begin{enumerate}
\conti
\item \texttt{savefig('sine.png')} saves the plot in,\pause
\begin{itemize}
\item The root directory ``/'' (on GNU/Linux, Unix based systems),
``C:'' (on windows).\pause
\item Will result in an error as full path is not supplied.\pause
\item The current working directory.\pause
\item Predefined directory like ``/documents''.
\end{itemize}
\end{enumerate}
\end{frame}
%%%%%%%%%%%%%%%%%%%%%%%%%%%%%%%%%%%%%%%%%%%%%%%%%%%%%%%%%%%%%%%%%%%%%%%%%%%%%%%%
\begin{frame}
\frametitle{Solutions}
\label{sec-11}
\begin{enumerate}
\item \texttt{savefig()}\pause
\item The current working directory
\end{enumerate}
\end{frame}
%%%%%%%%%%%%%%%%%%%%%%%%%%%%%%%%%%%%%%%%%%%%%%%%%%%%%%%%%%%%%%%%%%%%%%%%%%%%%%%%
\begin{frame}
\frametitle{FOSSEE}
{\color{blue}Free and Open-source Software for \\Science and Engineering Education} \\
\begin{itemize}
\item Goal - enabling all to use open source software tools
\item For more details, please visit {\color{blue}\url{http://fossee.in/}}
\end{itemize}
\end{frame}
%%%%%%%%%%%%%%%%%%%%%%%%%%%%%%%%%%%%%%%%%%%%%%%%%%%%%%%%%%%%%%%%%%%%%%%%%%%%%%%%
\begin{frame}
\frametitle{About the Spoken Tutorial Project}
\begin{itemize}
\item Watch the video available at {\color{blue}http://spoken-tutorial.org /What\_is\_a\_Spoken\_Tutorial}
\item It summarises the Spoken Tutorial project \pause
\item If you do not have good bandwidth, you can download and watch it
\end{itemize}
\end{frame}
%%%%%%%%%%%%%%%%%%%%%%%%%%%%%%%%%%%%%%%%%%%%%%%%%%%%%%%%%%%%%%%%%%%%%%%%%%%%%%%%
\begin{frame}
\frametitle{Spoken Tutorial Workshops}The Spoken Tutorial Project Team 
\begin{itemize}
\item Conducts workshops using spoken tutorials 
\item Gives certificates to those who pass an online test 
\item For more details, please write to \\ \hspace {0.5cm}{\color{blue}contact@spoken-tutorial.org}
\end{itemize}
\end{frame}
%%%%%%%%%%%%%%%%%%%%%%%%%%%%%%%%%%%%%%%%%%%%%%%%%%%%%%%%%%%%%%%%%%%%%%%%%%%%%%%%
\begin{frame}
\frametitle{Acknowledgements}
\begin{itemize}
\item Spoken Tutorial Project is a part of the Talk to a Teacher  project 
\item It is supported by the National Mission on Education through  ICT, MHRD, Government of India 
\item More information on this Mission is available at: \\{\color{blue}http://spoken-tutorial.org /NMEICT-Intro}
\end{itemize}
\end{frame}
%%%%%%%%%%%%%%%%%%%%%%%%%%%%%%%%%%%%%%%%%%%%%%%%%%%%%%%%%%%%%%%%%%%%%%%%%%%%%%%%
\begin{frame}

  \begin{block}{}
  \begin{center}
  \textcolor{blue}{\Large THANK YOU!} 
  \end{center}
  \end{block}
\begin{block}{}
  \begin{center}
    For more Information, visit our website\\
    \url{http://fossee.in/}
  \end{center}  
  \end{block}
\end{frame}
%%%%%%%%%%%%%%%%%%%%%%%%%%%%%%%%%%%%%%%%%%%%%%%%%%%%%%%%%%%%%%%%%%%%%%%%%%%%%%%%
\end{document}
